% ============================================================
% SECCIÓN 9: Conclusiones y Cierre de la Versión 1.0
% ============================================================
% Redactado en agosto de 2026 con los datos reales del cierre.
% Sustituye al esqueleto anterior, que asumía la ejecución del
% SPIKE-001 y 30 días de operación continua. Ninguna de las dos
% cosas llegó a ocurrir; el capítulo se reescribe en pasado.
% ============================================================

\chapter{Conclusiones y Cierre de la Versión 1.0}
\label{sec:conclusiones}

% ─────────────────────────────────────────────────────────────
\section{Estado de la Versión 1.0 Entregada}

Demeter v1.0 se entrega como un sistema \textbf{completo en su arquitectura y
verificado en banco}, modesto en alcance y nunca puesto en explotación real. Se
archiva como \textit{release} congelada: un artefacto estable, documentado y
reproducible, con una cláusula explícita de no continuidad.

\begin{table}[H]
    \centering
    \caption{Inventario factual de lo entregado en la versión 1.0.}
    \label{tab:inventario_v1}
    \begin{tabular}{lp{8cm}}
        \toprule
        \textbf{Artefacto} & \textbf{Estado al cierre} \\
        \midrule
        Capas implementadas & 4 de 4: firmware ESP32-S3, servicio de borde
        en Raspberry Pi, backend en contenedores y SDK científico \\
        Código propio & 27.025 líneas, excluidas dependencias y ficheros
        generados \\
        Protocolo & 14 tipos de mensaje operativos sobre trama binaria
        propia, con handshake y validación por capas \\
        Tests automatizados & 145, repartidos en cinco suites (véase el
        Capítulo~\ref{sec:testing}) \\
        Documentación técnica & Este documento, 84 páginas, 18 diagramas y
        9 cuadros de Decisión de Diseño \\
        Imágenes Docker & Backend y frontend publicadas en GitHub Container
        Registry \\
        Entornos de despliegue & 5 configuraciones de Docker Compose,
        levantables desde un árbol limpio \\
        Días de operación continua sin intervención & \textbf{0} \\
        \bottomrule
    \end{tabular}
\end{table}

Esa última fila es la que da la medida honesta del cierre y merece leerse junto a
todas las anteriores. El sistema \textbf{funciona}: la cadena completa, desde un
sensor que despierta de \textit{deep sleep} hasta el dato visible en el dashboard,
se recorrió de extremo a extremo sobre tres nodos reales. Pero funcionó en un banco
de laboratorio y bajo supervisión, nunca desatendido durante días ni expuesto a las
condiciones de un invernadero. Demeter v1.0 es un \textbf{prototipo validado}, no un
sistema en producción, y el resto de este capítulo se escribe desde esa premisa.

\begin{table}[H]
    \centering
    \caption{Alcance entregado frente a alcance excluido deliberadamente.}
    \label{tab:alcance_v1}
    \begin{tabular}{p{6.2cm}p{6.2cm}}
        \toprule
        \textbf{Entregado en v1.0} & \textbf{Fuera de alcance por decisión} \\
        \midrule
        Telemetría ambiental y de suelo por nodo y por planta &
        Cifrado del enlace de campo (AES-128 sobre ESP-NOW) \\
        Actuación remota sobre GPIO y secuencias programadas &
        Actualización de firmware por aire (OTA) \\
        Trazabilidad LIMS de plantas y experimentos &
        Multi-invernadero sobre una misma instancia \\
        SDK con cálculo agronómico y exportación &
        Observabilidad centralizada y sistema de alarmas \\
        Despliegue reproducible en cinco entornos &
        Interfaz adaptada a móvil \\
        \bottomrule
    \end{tabular}
\end{table}

% ─────────────────────────────────────────────────────────────
\section{Cumplimiento de Objetivos}

El Capítulo~\ref{sec:introduccion} declaró cuatro objetivos específicos. Su estado
real al cierre es el siguiente.

\begin{table}[H]
    \centering
    \caption{Grado de cumplimiento de los objetivos declarados.}
    \label{tab:cumplimiento}
    \begin{tabular}{p{2.8cm}p{2.6cm}p{7cm}}
        \toprule
        \textbf{Objetivo} & \textbf{Estado} & \textbf{Evidencia y matices} \\
        \midrule
        Automatización IoT & Cumplido parcialmente &
        Tres nodos ESP32-S3 operativos en banco, con ciclo de \textit{deep sleep}
        y reporte automático. La recolección «continua y sin intervención humana»
        que enunciaba el objetivo nunca se sostuvo más allá de una sesión de
        pruebas. \\
        \midrule
        Control Remoto & Cumplido parcialmente &
        La cadena de mando completa opera para \texttt{SET\_GPIO}, y el
        secuenciador ejecuta secuencias multi-paso. Pero \textbf{no se actuó nunca
        sobre una bomba real}: la instalación hidráulica se adquirió y no llegó a
        montarse, y el \textit{spike} que debía validarla no se ejecutó. Este es
        el objetivo que queda más lejos de lo prometido. \\
        \midrule
        Trazabilidad LIMS & Cumplido &
        Modelo de 13 tablas con relación N:M entre plantas y experimentos,
        consulta de telemetría por planta y registro de actividad
        (Capítulo~\ref{sec:backend}). \\
        \midrule
        Análisis Científico & Cumplido &
        SDK con caché Parquet, pipeline ETL, cálculo agronómico vectorizado y
        exportación a CSV y Excel (Capítulo~\ref{sec:sdk}), utilizable sin
        conocimientos de la infraestructura. \\
        \bottomrule
    \end{tabular}
\end{table}

El patrón es nítido y conviene nombrarlo: \textbf{los dos objetivos de software puro
se cumplieron; los dos que tocaban el mundo físico se quedaron a medias}. No es
casualidad. El software se puede escribir de noche, en solitario y sin depender de
nadie; montar una instalación de riego en un invernadero de otra escuela exige
coordinación, permisos, presencia física y manos ajenas. La barrera no fue técnica.

% ─────────────────────────────────────────────────────────────
\section{Decisiones Arquitectónicas en Retrospectiva}

\subsection*{Decisiones que se sostienen}

\begin{itemize}[leftmargin=*, label=--]
    \item \textbf{Protocolo binario propio en lugar de JSON sobre el enlace de
          campo.} La trama compacta hizo viable el \textit{deep sleep} agresivo y
          mantuvo la ocupación del canal por nodo por debajo del 0,01\,\%. El coste
          --- tener que escribir y depurar el empaquetado --- se amortizó con la
          validación por capas, que atrapa tramas corruptas antes de que lleguen a
          la lógica de dominio.
    \item \textbf{Pydantic como puente único entre lo binario y el dominio.}
          Definir el contrato una sola vez en \texttt{Software/Common} y consumirlo
          desde el backend y desde la Raspberry Pi eliminó de raíz toda una clase
          de errores de desincronización entre extremos. Es, con diferencia, la
          decisión que más problemas evitó.
    \item \textbf{Patrón \textit{Strategy} sobre la interfaz \texttt{IComms}.}
          Permitió probar la lógica de protocolo en el PC, sin hardware, y explica
          que la suite de firmware sea la segunda más poblada del proyecto.
    \item \textbf{Imagen base común para nube y Raspberry Pi.} La misma
          definición de servicios sirvió para cinco entornos, con las diferencias
          confinadas a ficheros de \textit{override}.
    \item \textbf{Caché Parquet local en el SDK.} Hace que el análisis iterativo
          no golpee la API en cada ejecución, que es exactamente el patrón de
          trabajo de quien explora datos.
\end{itemize}

\subsection*{Decisiones que se tomarían de otra forma}

\begin{itemize}[leftmargin=*, label=--]
    \item \textbf{Validar el hardware crítico antes de comprarlo.} El
          \textit{spike} de la bomba se escribió y nunca se ejecutó, y para cuando
          se adquirió la instalación hidráulica su plan de pruebas ya estaba
          obsoleto: contemplaba una bomba de 5\,V y el material comprado era de
          12\,V. Un \textit{spike} que no se ejecuta no reduce ningún riesgo; solo
          documenta que se era consciente de él.
    \item \textbf{Medir desde el primer día.} El sistema se dio por bueno sin
          línea base de latencia, pérdida de tramas ni autonomía real. Sin esas
          medidas no hay forma de demostrar que funciona ni de detectar una
          regresión, y añadirlas al final ya no es posible.
    \item \textbf{Integración continua completa desde el primer
          \textit{release}.} El \textit{pipeline} construye y publica las imágenes
          de backend y frontend, pero nunca compiló el firmware ni ejecutó
          análisis de secretos. La auditoría del cierre encontró una credencial
          viva versionada en el repositorio, que un análisis automático habría
          detectado el día que se introdujo.
    \item \textbf{Higiene del repositorio desde el inicio.} Versionar
          \texttt{node\_modules}, la documentación generada por Doxygen y ficheros
          de \textit{log} infla el repositorio en un orden de magnitud y arruina
          cualquier métrica automática sobre él. Corregirlo después obliga a
          reescribir el histórico.
\end{itemize}

\subsection*{Sobre las decisiones de arquitectura y su registro}

El plan original contemplaba formalizar entre cinco y siete \textit{Architecture
Decision Records} como documentos independientes. \textbf{No se escribió ninguno.}
En la práctica su función la cumplen los \textbf{nueve cuadros de Decisión de
Diseño} repartidos por este documento, que recogen el mismo contenido --- decisión,
alternativas descartadas y justificación --- pero incrustado en el capítulo donde
la decisión es relevante en lugar de en un registro aparte.

La lección no es que los ADR sobren, sino que \textbf{un formato que nadie mantiene
es peor que uno más humilde que sí se usa}. Para un proyecto de un solo autor, el
cuadro dentro del capítulo resultó ser el vehículo que sobrevivió; el registro
separado se habría quedado, como tantas otras cosas, en la lista de pendientes.

% ─────────────────────────────────────────────────────────────
\section{Lecciones Aprendidas}

\subsection*{Técnicas}

\begin{itemize}[leftmargin=*, label=--]
    \item \textbf{El contrato compartido vale más que cualquier optimización.}
          Un único modelo de datos consumido por todas las capas evitó la clase de
          fallo más cara de depurar en un sistema distribuido: dos extremos que
          creen hablar el mismo idioma y no lo hacen.
    \item \textbf{\texttt{asyncio} recompensa la disciplina y castiga la
          improvisación.} Cuatro tareas cooperativas sobre un solo hilo bastaron
          para el servicio de borde, pero cualquier llamada bloqueante que se cuele
          detiene todo el proceso. El modelo obliga a pensar dónde está cada
          punto de espera.
    \item \textbf{ESP-NOW resuelve el enlace de campo mejor de lo esperado, y con
          un límite claro.} No necesita punto de acceso, el consumo es compatible
          con el \textit{deep sleep} y el alcance fue suficiente. Su límite es la
          seguridad: sin cifrado, cualquier receptor cercano lee la trama.
    \item \textbf{El binario merece la pena solo cuando la energía es el
          recurso escaso.} En el enlace de campo, sí. Entre la Raspberry Pi y el
          servidor, donde ya había JSON sobre WebSocket, habría sido una
          complicación gratuita.
\end{itemize}

\subsection*{De proceso}

\begin{itemize}[leftmargin=*, label=--]
    \item \textbf{Documentación como código.} Mantener el LaTeX en el mismo
          repositorio, versionado junto al código que describe, es lo que ha hecho
          posible que este documento siga siendo cierto en el momento del cierre.
          Una documentación fuera del repositorio habría divergido en semanas.
    \item \textbf{Un \textit{spike} sin ejecutar es una deuda disfrazada de
          diligencia.} Escribir la exploración y no hacerla da la sensación de
          haber gestionado el riesgo sin haberlo tocado. Es peor que no escribirla,
          porque tranquiliza.
    \item \textbf{Cerrar es una decisión de ingeniería, no una rendición.} La
          alternativa realista a este cierre no era un proyecto que siguiera
          creciendo, sino uno que siguiera facturando infraestructura sin que nadie
          lo tocara. Decidir el final y documentarlo produce un artefacto
          referenciable; dejarlo morir por inercia no produce nada.
\end{itemize}

\subsection*{Personales}

El aprendizaje más incómodo del proyecto no es técnico. Demeter demostró que una
sola persona puede construir un sistema distribuido de cuatro capas, documentarlo
en ochenta páginas y desplegarlo de forma reproducible. Demostró también que \textbf{esa
misma persona no puede sostenerlo}, y que la parte que exigía coordinarse con
otros --- montar la instalación en el invernadero, medir en campo, transferir el
conocimiento --- fue precisamente la que no se hizo.

La conclusión que el autor se lleva al siguiente proyecto es que la capacidad
técnica individual no es el cuello de botella, y que optimizarla más no habría
cambiado el resultado. Lo que faltó fue todo lo demás.

% ─────────────────────────────────────────────────────────────
\section{Trabajo Futuro y Extensiones Posibles}

Las siguientes extensiones son plausibles sobre la arquitectura entregada. Se
inventarían aquí para quien retome el trabajo; \textbf{ninguna constituye un
compromiso} de la versión 1.0 ni del autor.

\begin{itemize}[leftmargin=*, label=--]
    \item \textbf{Caracterización en campo.} Latencia extremo a extremo, pérdida
          de tramas, autonomía real y tiempo de reconexión. Es el primer trabajo
          que debería hacer cualquier continuación: sin línea base no se puede
          saber si un cambio mejora o empeora el sistema.
    \item \textbf{Montaje y validación de la instalación de riego.} El material
          está adquirido y sin usar. Requiere reescribir el \textit{spike} para
          12\,V y dimensionar la conmutación: la fuente disponible no admite más de
          una bomba simultánea.
    \item \textbf{Cifrado del enlace de campo.} AES-128 sobre el \textit{payload}
          ESP-NOW, imprescindible antes de cualquier despliegue con datos que
          importen.
    \item \textbf{Observabilidad centralizada.} Agregación de \textit{logs} y
          métricas. Hoy diagnosticar una incidencia obliga a leer la salida de
          cada contenedor por separado.
    \item \textbf{Alta de nodos sin recompilar.} La dirección del \textit{gateway}
          y el identificador de nodo se fijan en tiempo de compilación; añadir un
          nodo exige reflashear.
    \item \textbf{Actualización de firmware por aire.} Con autenticación y
          \textit{rollback}. Se descartó en v1.0 por alcance, pero condiciona
          cualquier despliegue de más de un puñado de nodos.
    \item \textbf{Riego predictivo.} Integración con datos meteorológicos o una
          capa de inferencia por especie. Es la extensión más vistosa y la que
          menos sentido tiene abordar antes que las cinco anteriores.
\end{itemize}

% ─────────────────────────────────────────────────────────────
\section{Cierre y Agradecimientos}

Con la publicación de este documento, el Proyecto Demeter queda \textbf{archivado
en su versión 1.0}. No hay desarrollo planificado. El repositorio se abre con
licencia MIT para el código y CC-BY-SA 4.0 para la documentación, de modo que
cualquiera pueda estudiarlo, reutilizarlo o continuarlo por su cuenta, sin que ello
comprometa al autor ni a la asociación a mantenimiento alguno.

El análisis que sostiene esta decisión --- económico, de recursos humanos y técnico
--- se recoge en el documento hermano \textit{Informe de Viabilidad y Propuesta de
Cierre}, en \texttt{docs/Closure/}.

Quedan los agradecimientos. A la \textbf{Asociación Universitaria ESIBot}, por el
contexto, los recursos y la libertad para llevar el proyecto hasta donde llegó. Al
\textbf{Departamento de Agronomía de la ETSIA}, por plantear un problema real y
confiar en que un grupo de estudiantes de otra escuela pudiera abordarlo. Y a
quienes, en algún momento de estos veinte meses, preguntaron cómo iba: sostener un
proyecto en solitario se parece bastante a hablar solo, y esas preguntas contaban
más de lo que probablemente parecía.

\begin{designnote}
\textbf{Por qué se archiva como \textit{release} y no como abandono.} La diferencia
es material. Un proyecto abandonado deja de ser operable: nadie sabe si arranca, la
documentación miente y el conocimiento se disipa. Una versión archivada es un
artefacto estable, reproducible desde un árbol limpio, documentado hasta sus
limitaciones y publicado con licencia clara. Demeter cumple los criterios de la
segunda categoría, incluido el más incómodo: dejar escrito, con nombres y cifras,
todo aquello que no llegó a funcionar. Un documento que solo contara los aciertos
sería propaganda; y la propaganda, a diferencia de un artefacto, no le sirve a
nadie que venga después.
\end{designnote}
